\label{subsec:zeta-online-kernel-part2}

such that for all invertible $r$ the identity
\[
\widehat{\Delta}(w,r+r^{-1})=\Delta(w/r,r^2).
\]
holds. Let $\Phi_m^+(s)$ denote the minimal polynomial (with integer coefficients, monic) of the algebraic number $s_m:=2\cos(\pi/m)$, and define
\[
R_m(w):=\operatorname{Res}_s\bigl(\widehat{\Delta}(w,s),\Phi_m^+(s)\bigr)\in\ZZ[w].
\]
Then when $u=e^{2\pi i/m}$, the solutions $z$ of $\Delta(z,u)=0$ (equivalently $w=z\sqrt{u}$) must satisfy $R_m(w)=0$. Moreover, under the unit-circle twist $|u|=1$,
\[
\rho_m=\frac{1}{\min\{|w|:R_m(w)=0\}}.
\]
\end{proposition}

\begin{remark}[Invariant subring of the completion: the multiplicative twist \texorpdfstring{$u$}{u} is compressed to the additive coordinate \texorpdfstring{$s$}{s}]\label{rem:completion-invariant-ring}
At the level of variables, the substitution used in the completion
\(
u=r^2,\ s=r+r^{-1}
\)
can be understood as invariantization under the involution \(r\mapsto r^{-1}\): the fixed subring of the Laurent ring satisfies
\[
\ZZ[r,r^{-1}]^{\,r\mapsto r^{-1}}
=\ZZ\!\left[r+r^{-1}\right]
=\ZZ[s].
\]
Therefore, under the self-dual symmetry (\(u\leftrightarrow 1/u\)), the multiplicative twist parameter \(u\in\CC^\times\) at a root of unity is compressed via completion into a real ``trace coordinate'' \(s=2\cos(\pi/m)\), and \(\widehat\Delta(w,s)\in\ZZ[s][w]\) provides a unified interface for transcribing the ``multiplicative-side twisted spectrum'' into ``additive-side algebraic certificates.''
\end{remark}

\begin{remark}[Explicit formula for the completed determinant of the synchronization kernel and the \texorpdfstring{$s\mapsto -s$}{s->-s} symmetry]\label{rem:sync-hatdelta-explicit}
For the synchronization kernel with the output bit potential weighting, Appendix \ref{app:pressure-analytic} provides the characteristic polynomial of the nonzero spectrum $F(\lambda,u)\in\ZZ[u][\lambda]$, and
\[
\Delta(z,u)=z^6F(1/z,u)
\]
(equivalently, $B(u)$ has only $6$ nonzero eigenvalues). Taking the completed variables $u=r^2,\ w=zr,\ s=r+r^{-1}$, the completed determinant of Proposition \ref{prop:sync-cyclotomic-elimination} can be written as the following \textbf{integer-coefficient closed form}:
\[
\boxed{\;
\widehat\Delta(w,s)=
1-sw-5w^2+3sw^3+(5-s^2)w^4+(s^3-6s)w^5+(s^2-1)w^6.
\;}
\]
Moreover, direct substitution verifies that it satisfies
\[
\widehat\Delta(w,-s)=\widehat\Delta(-w,s).
\]
Therefore, when $s$ changes sign, the $w$-root set corresponds under $w\mapsto -w$, and the moduli of the zeros are unchanged; consequently, the ``minimum modulus zero'' defined by $\min\{|w|:\widehat\Delta(w,s)=0\}$ depends only on $|s|$.
\end{remark}

\begin{corollary}[Closed-form factorization and minimum modulus root at the endpoint \(u=1\) (\(s=2\))]\label{cor:sync-hatdelta-s2-factorization}
Setting \(s=2\), one obtains
\[
\widehat\Delta(w,2)
=1-2w-5w^2+6w^3+w^4-4w^5+3w^6
=(w-1)(w+1)(3w-1)(w^3-w^2+w+1).
\]
Therefore, among the roots of the equation \(\widehat\Delta(w,2)=0\), the one with minimum \(|w|\) is
\[
\boxed{\ w_\star(2)=\frac13\ },\qquad \rho(2)=\frac{1}{|w_\star(2)|}=3.
\]
\end{corollary}
\begin{proof}
Direct substitution and factorization yield the stated decomposition.
To verify that the minimum modulus root is indeed \(1/3\), note that when \(|w|<1/3\),
\[
\bigl|w^3-w^2+w\bigr|
\le |w|+|w|^2+|w|^3
<\frac13+\frac19+\frac1{27}
=\frac{13}{27}<1,
\]
so \(w^3-w^2+w+1\neq 0\). Since the zeros of \((w\pm 1)=0\) have modulus \(1\), the minimum possible value of \(|w|\) can only come from the factor \((3w-1)\), i.e., \(w_\star(2)=1/3\).
Thus \(\rho(2)=1/|w_\star(2)|=3\). \qed
\end{proof}

\begin{corollary}[Closed-form spectral radius for the worst-case twist \(m=2\)]\label{cor:sync-rho-m2-closed-form}
Setting \(m=2\), one has \(u=e^{2\pi i/m}=-1\) and
\(
s_2=2\cos(\pi/2)=0
\).
Then
\[
R_2(w)=\widehat\Delta(w,0)=1-5w^2+5w^4-w^6=(1-w^2)(w^4-4w^2+1).
\]
Let \(t:=w^2\); then \(R_2(w)=0\) is equivalent to
\[
(1-t)(t^2-4t+1)=0,
\]
whose positive roots are \(t=1\) and \(t=2\pm\sqrt3\). Therefore, the minimum modulus zero satisfies
\[
\min\{|w|:R_2(w)=0\}=\sqrt{2-\sqrt3},
\]
yielding the closed-form constant
\[
\boxed{\ \rho_2=\frac{1}{\sqrt{2-\sqrt3}}=\sqrt{2+\sqrt3}\approx 1.9318516526\ }.
\]
\end{corollary}
\begin{proof}
By substituting \(s=0\) into the explicit formula from Remark \ref{rem:sync-hatdelta-explicit} and factoring, one obtains the stated \(R_2(w)\).
By definition
\(
\rho_2=1/\min\{|w|:R_2(w)=0\}
\),
and the conclusion follows immediately from the explicit root computation. \qed
\end{proof}

\begin{remark}[Endpoint derivative fingerprint: \texorpdfstring{$\rho'(2)=11/17$}{rho'(2)=11/17}]\label{rem:sync-rho-derivative-fingerprint}
Let \(w_\star(s)\) denote the root of the equation \(\widehat\Delta(w,s)=0\) with minimum modulus in the \(w\)-plane (corresponding to the spectral radius ratio of the unit-circle twist channel \(\rho(s)=1/|w_\star(s)|\)).
At the endpoint \(s=2\), direct substitution yields \(w_\star(2)=1/3\), and this root is simple.
By the implicit function theorem, there exists a unique analytic branch \(w(s)\) (in some neighborhood of \(s=2\)) satisfying
\(\widehat\Delta(w(s),s)=0\) and \(w(2)=1/3\);
setting \(\rho(s):=1/w(s)\), explicit implicit differentiation yields
\[
w'(2)=-\frac{\partial_s\widehat\Delta}{\partial_w\widehat\Delta}\Big|_{(w,s)=(1/3,\,2)}=-\frac{11}{153},
\qquad
\boxed{\ \rho'(2)=\frac{11}{17}\ }.
\]
Further second-order differentiation gives
\[
\boxed{\ \rho''(2)=\frac{1351}{9826}\ }.
\]
Further third-order differentiation (or implicit function expansion at \((w,s)=(1/3,2)\)) yields the local Taylor expansion of the ``dispersion curve'': letting \(\delta:=2-s\),
\[
\boxed{\;
\rho(2-\delta)=3-\frac{11}{17}\delta+\frac{1351}{19652}\delta^2-\frac{17133}{45435424}\delta^3+O(\delta^4).\;}
\]
Combined with the endpoint approximation
\(
s_m=2\cos(\pi/m)=2-\pi^2/m^2+\pi^4/(12m^4)+O(m^{-6})
\),
Taylor expanding \(\rho(s_m)\) recovers
\[
\rho_m
=3-\frac{11\pi^2}{17m^2}
+\frac{1808\pi^4}{14739m^4}
-\frac{27872249\pi^6}{2044594080m^6}
+O(m^{-8}),
\]
consistent with the constants in Corollary \ref{cor:sync-rho-gap-m2} and Corollary \ref{cor:sync-rho-gap-m8-diffusive}.
\end{remark}

\begin{proposition}[Diffusion dictionary of endpoint derivatives: \texorpdfstring{$\Lambda'(2)$}{Lambda'(2)} locks the quadratic dissipation of the unit-circle twist and the near-\texorpdfstring{$1$}{1} dominant singularity trajectory]\label{prop:completion-endpoint-jet-diffusion}
Assume that a certain completion/unit-circle twist channel satisfies strict phase--amplitude separation: there exists a real function \(g(s)=\log\Lambda(s)\) analytic in some neighborhood of \(s=2\) (with \(\Lambda(2)>1\)) such that for \(u=e^{it}\),
\[
P(it)=\log\lambda(e^{it})=\frac{it}{2}+g\!\left(S(t)\right),
\qquad
S(t):=2\cos(t/2).
\]
Denote \(\lambda:=\Lambda(2)\), and define the unit-circle spectral ratio
\[
\eta(t):=\frac{|\lambda(e^{it})|}{\lambda}
=\exp\!\bigl(\Re P(it)-\log\lambda\bigr)
=\frac{\Lambda(S(t))}{\Lambda(2)}.
\]
Then as \(t\to 0\), the endpoint jet expansion reads
\[
\boxed{\;
\log\eta(t)
=-\frac{g'(2)}{4}\,t^{2}
+\left(\frac{g'(2)}{192}+\frac{g''(2)}{32}\right)t^{4}
+O(t^{6})\;}
\]
so that the quadratic ``diffusion variance'' and dissipation constant are given in closed form by the endpoint derivatives:
\[
\boxed{\;
\Re P(it)=\log\lambda-\frac{\sigma^{2}}{2}t^{2}+O(t^{4}),\qquad
\sigma^{2}=\frac{g'(2)}{2}=\frac{\Lambda'(2)}{2\Lambda(2)},\qquad
\kappa_\infty:=\frac{\sigma^{2}}{2}=\frac{\Lambda'(2)}{4\Lambda(2)}\; }.
\]
Further define the normalized trajectory pulling back the unit-circle twist's dominant singularity to the \(s\)-plane:
\[
s_{\mathrm{pole}}(t):=\frac{P(it)}{\log\lambda},
\]
then
\[
\boxed{\;
s_{\mathrm{pole}}(t)
=1+i\,\frac{t}{2\log\lambda}-\frac{\kappa_\infty}{\log\lambda}\,t^{2}+O(t^{4})\;}
\]
where the drift coefficient \(1/2\) is the universal output of phase--amplitude separation, and the dissipation term is entirely determined by the endpoint derivative \(\Lambda'(2)\).
\end{proposition}
\begin{proof}
From \(S(t)=2\cos(t/2)=2-\frac{t^{2}}{4}+\frac{t^{4}}{192}+O(t^{6})\) and the Taylor expansion of \(g\) at \(2\), the result follows.
Since \(|e^{it/2}|=1\), one has \(\Re P(it)=g(S(t))\), whence \(\eta(t)=\exp(g(S(t))-g(2))\).
Finally, \(s_{\mathrm{pole}}(t)=1+\frac{it}{2\log\lambda}+\frac{g(S(t))-g(2)}{\log\lambda}\), and substituting the above expansion completes the proof. \qed
\end{proof}

\begin{remark}[Self-check: the synchronization kernel endpoint derivative \texorpdfstring{$\rho'(2)$}{rho'(2)} immediately yields \texorpdfstring{$11/204$}{11/204}]\label{rem:sync-endpoint-diffusion-check}
In the synchronization kernel case, one can take \(\Lambda(s)=\rho(s)\), \(\lambda=\rho(2)=3\), and Remark \ref{rem:sync-rho-derivative-fingerprint} gives \(\rho'(2)=11/17\).
By Proposition \ref{prop:completion-endpoint-jet-diffusion},
\[
\kappa_\infty=\frac{\rho'(2)}{4\rho(2)}=\frac{11/17}{12}=\frac{11}{204},
\]
which is fully consistent with the quadratic coefficient in Remark \ref{rem:sync-rho-pressure-asymp} and Corollary \ref{cor:sync-rho-gap-m2}.
\end{remark}

\begin{remark}[Algebraic certificates for spectral bifurcation points: real roots of \texorpdfstring{$\mathrm{Disc}_w(\widehat\Delta)$}{Disc}]\label{rem:sync-hatdelta-discriminant-branch}
For the degree-six polynomial in \(w\), \(\widehat\Delta(w,s)\), its discriminant \(\mathrm{Disc}_w(\widehat\Delta)(s)\in\ZZ[s]\) has zeros that characterize the parameter values \(s\) at which ``repeated roots exist,'' i.e., the candidate locations where spectral branches collide/exchange in the \(s\)-direction.
The following formula and the corresponding positive real roots (restricted to \(s\in(0,2]\), corresponding to unit-circle twist \(s=2\cos(\theta/2)\)) are exported in a single step by the script \texttt{scripts/exp\_sync\_kernel\_hatdelta\_discriminant.py}:
\input{sections/generated/eq_sync_kernel_hatdelta_discriminant}
\input{sections/generated/tab_sync_kernel_hatdelta_branch_points}
\end{remark}

\begin{remark}[True analytic radius: the nearest complex bifurcation point pins down the convergence disk at \texorpdfstring{$t=0$}{t=0}]\label{rem:sync-hatdelta-analytic-radius-complex}
Table \ref{tab:sync_kernel_hatdelta_branch_points} provides the three bifurcation angles of \(\mathrm{Disc}_w(\widehat\Delta)(s)=0\) on the real axis \(s\in(0,2]\).
To avoid confusion with the variable \(\theta\) in the pressure expansion of Remark \ref{rem:sync-rho-pressure-asymp}, this remark uses \(t\) to denote the unit-circle twist angle (\(u=e^{it}\), i.e., the angular coordinate \(\theta\) in Table \ref{tab:sync_kernel_hatdelta_branch_points}).
For any local expansion around \(t=0\) (including the small-angle continuation with \(\theta=it\) in Remark \ref{rem:sync-rho-pressure-asymp}), the true convergence radius is determined by the nearest analytic obstruction in the \emph{complex domain}.
\par
Since \eqref{eq:sync_kernel_hatdelta_discriminant} contains only even powers of \(s\), one can set \(x:=s^2\) to reduce the discriminant to a degree-\(10\) polynomial and find its roots in the complex plane; pulling back the nearest pair of conjugate roots via
\(
s=\sqrt{x},\
t=2\arccos(s/2)
\)
(principal branch) to the angular plane, one obtains the nearest complex bifurcation point \(t_\star,\overline{t_\star}\) to \(t=0\) and the true analytic radius
\[
\boxed{\;R_t:=|t_\star|\approx 1.070953\;}
\]
(see Table \ref{tab:sync_kernel_hatdelta_nearest_complex_branch_point}; the angular coordinate is still denoted \(\theta_\star\equiv t_\star\)).
Note that \(R_t\) is strictly smaller than the smallest real angle in Table \ref{tab:sync_kernel_hatdelta_branch_points} (\(\approx 2.03063\)); therefore, any Taylor/cumulant/continuation expansion centered at \(t=0\) has its convergence domain first truncated by this complex bifurcation point, rather than by the smallest real bifurcation angle.
\par
This yields a completely internal, auditable ``modulus selection law'': when the root-of-unity twist angle is \(t_m:=2\pi/m\),
\[
t_m<R_t\ \Longrightarrow\ \text{the local expansion converges strictly}\quad(\text{in particular }m\ge 6),
\]
where at \(m=6\) one has \(t_6/R_t\approx 0.978\), which is close to the convergence boundary and causes the local expansion to converge slowly, while \(m=5\) already falls outside the convergence disk.
Moreover, the real part of the nearest complex bifurcation point in the \(s\)-plane differs from \(s_6=2\cos(\pi/6)=\sqrt3\) by a very small amount (Table \ref{tab:sync_kernel_hatdelta_nearest_complex_branch_point} gives \(\Re(s_\star)-\sqrt3\approx 2.76\times 10^{-3}\)); this tiny deviation is precisely determined by the complex coordinate of the bifurcation point, thereby numerically explaining the criticality of the convergence properties at \(m=6\).
\par
Therefore, when computing \(\rho_m\) (or the residue/Dirichlet--Mertens constants derived from it), the choice between ``local expansion'' and ``global stripping'' can be written as an auditable strategy fork:
\begin{itemize}
  \item \textbf{\(m\le 5\):} \(t_m>R_t\), the local power series does not converge; one should proceed directly via the global path (e.g., the cyclotomic resultant root-finding/stripping of Remark \ref{rem:sync-cyclotomic-elimination-table}).
  \item \textbf{\(m=6\):} \(t_6/R_t\approx 0.978\) is close to the convergence boundary; the local expansion, while within the convergence disk, converges slowly, and should be combined with acceleration (Pad\'e, branch stripping) or directly fall back to global root-finding.
  \item \textbf{\(m\ge 7\):} enters the robust regime; the dominant factor of the local truncation error can be budgeted as the geometric ratio \((t_m/R_t)^K\) (see the small-angle prediction interface of Remark \ref{rem:sync-rho-pressure-asymp}).
 \end{itemize}
\input{sections/generated/tab_sync_kernel_hatdelta_nearest_complex_branch_point}
\par
Further, using the completed log variable \(\theta:=\log u=it\) as the analytic continuation coordinate and introducing the half-angle compression coordinate
\(
x:=\tan(\theta/2)
\).
The nearest complex bifurcation point in the \(\theta\)-plane is then mapped to a pair of conjugate singularities \(x_\star,\overline{x_\star}\) in the \(x\)-plane; their modulus \(|x_\star|\approx 0.4916\) gives the maximum analytic disk radius with \(x\) as the local coordinate (thereby providing an auditable convergence domain guardrail for Taylor/Pad\'e organization based on \(x\)).
\input{sections/generated/tab_sync_kernel_hatdelta_nearest_complex_branch_point_tan_half_angle}
\end{remark}

\begin{remark}[Reproducible output: \texorpdfstring{$R_m$}{Rm} and \texorpdfstring{$\rho_m$}{rhom} for \texorpdfstring{$m=3,\dots,12$}{m=3..12}]\label{rem:sync-cyclotomic-elimination-table}
The script \texttt{scripts/exp\_sync\_kernel\_cyclotomic\_elimination.py} constructs $\widehat{\Delta}(w,s)$ from the $F(\lambda,u)$ of Appendix \ref{app:pressure-analytic}, and computes $R_m(w)$ and $\rho_m$ for $m=3,\dots,12$; the output is given in Table \ref{tab:sync_kernel_cyclotomic_elimination_summary} and the polynomial block \texttt{tab\_sync\_kernel\_cyclotomic\_elimination\_polys.tex} (which can be used directly for regression testing).
\end{remark}

\begin{table}[H]
\centering
\scriptsize
\setlength{\tabcolsep}{6pt}
\caption{Cyclotomic elimination for mod-$m$ prime-shadow mixing exponents of the weighted sync-kernel. $R_m(w)$ is an integer resultant polynomial; $\rho_m$ is obtained as $\rho_m=1/\min\{|w|:R_m(w)=0\}$. We also list the pressure-based prediction $\rho_m^{\mathrm{pred}}$ from the Taylor expansion of $P(\theta)$.}
\label{tab:sync_kernel_cyclotomic_elimination_summary}
\begin{tabular}{r r r r r}
\toprule
$m$ & $\deg R_m$ & $\rho_m$ (resultant) & $\rho_m^{\mathrm{pred}}$ & abs err\\
\midrule
3 & 5 & 2.418034345100 & 2.418835165986 & 8.008e-04\\
4 & 12 & 2.644080451866 & 2.643950096830 & 1.304e-04\\
5 & 12 & 2.762782720637 & 2.762705118214 & 7.760e-05\\
6 & 12 & 2.831532258235 & 2.831497400579 & 3.486e-05\\
7 & 18 & 2.874530783058 & 2.874514708145 & 1.607e-05\\
8 & 24 & 2.903081377378 & 2.903073483799 & 7.894e-06\\
9 & 18 & 2.922953899166 & 2.922949770585 & 4.129e-06\\
10 & 24 & 2.937319429068 & 2.937317144301 & 2.285e-06\\
11 & 30 & 2.948029997898 & 2.948028669926 & 1.328e-06\\
12 & 24 & 2.956223089850 & 2.956222284562 & 8.053e-07\\
\bottomrule
\end{tabular}
\end{table}

% AUTO-GENERATED by scripts/exp_sync_kernel_cyclotomic_elimination.py
\begin{align*}
R_{3}(w)&=5 w^{5} - 4 w^{4} - 3 w^{3} + 5 w^{2} + w - 1\\
R_{4}(w)&=w^{12} - 26 w^{10} + 47 w^{8} - 62 w^{6} + 43 w^{4} - 12 w^{2} + 1\\
R_{5}(w)&=w^{12} - 9 w^{11} + 13 w^{10} + 23 w^{9} - 33 w^{8} - 30 w^{7} + 52 w^{6} + 23 w^{5} - 38 w^{4} - 8 w^{3} + 11 w^{2} + w - 1\\
R_{6}(w)&=4 w^{12} - 19 w^{10} + 38 w^{8} - 61 w^{6} + 47 w^{4} - 13 w^{2} + 1\\
R_{7}(w)&=w^{18} + 11 w^{17} + 4 w^{16} - 107 w^{15} + 83 w^{14} + 249 w^{13} - 288 w^{12} - 352 w^{11} + 492 w^{10} + 338 w^{9} - 526 w^{8} - 208 w^{7} + 325 w^{6} + 75 w^{5} - 107 w^{4} - 14 w^{3} + 17 w^{2} + w - 1\\
\end{align*}


\begin{remark}[The pressure expansion yields high-precision asymptotic predictions for \texorpdfstring{$\rho_m$}{rhom}]\label{rem:sync-rho-pressure-asymp}
From the Taylor expansion in Appendix \ref{app:pressure-analytic},
\[
\begin{aligned}
P(\theta)=\log 3+\frac12\theta+\frac{11}{204}\theta^2+\frac{1559}{1414944}\theta^4\\
-\frac{17123893}{1177686190080}\theta^6-\frac{122803509253}{25412897733672960}\theta^8\\
-\frac{518906628614669}{10575831520845339033600}\theta^{10}
+\frac{24353138488976295223}{880247609142999258444595200}\theta^{12}
+O(\theta^{14}),
\end{aligned}
\]
setting $\theta=it$ and writing $u=e^{it}$, one obtains
\[
\begin{aligned}
\log|\lambda(e^{it})|=\Re P(it)=\log 3-\frac{11}{204}t^2+\frac{1559}{1414944}t^4\\
+\frac{17123893}{1177686190080}t^6-\frac{122803509253}{25412897733672960}t^8\\
+\frac{518906628614669}{10575831520845339033600}t^{10}
+\frac{24353138488976295223}{880247609142999258444595200}t^{12}
+O(t^{14}).
\end{aligned}
\]
Setting $t=2\pi/m$ gives the directly usable prediction formula
\[
\begin{aligned}
\rho_m\approx 3\exp\!\Bigl(
&-\frac{11}{204}\Bigl(\frac{2\pi}{m}\Bigr)^2+\frac{1559}{1414944}\Bigl(\frac{2\pi}{m}\Bigr)^4\\
&+\frac{17123893}{1177686190080}\Bigl(\frac{2\pi}{m}\Bigr)^6-\frac{122803509253}{25412897733672960}\Bigl(\frac{2\pi}{m}\Bigr)^8\\
&+\frac{518906628614669}{10575831520845339033600}\Bigl(\frac{2\pi}{m}\Bigr)^{10}
+\frac{24353138488976295223}{880247609142999258444595200}\Bigl(\frac{2\pi}{m}\Bigr)^{12}
\Bigr),
\end{aligned}
\]
whose agreement with the exact resultant values above is extremely high: for example, defining the ``true value'' $\rho_m$ via the minimum modulus zero of $\widehat\Delta(w,s_m)=0$, one has
\[
\frac{|\rho_5^{\mathrm{pred}}-\rho_5|}{\rho_5}\approx 3.99\times 10^{-7},\qquad
\frac{|\rho_{10}^{\mathrm{pred}}-\rho_{10}|}{\rho_{10}}\approx 1.03\times 10^{-10},\qquad
\frac{|\rho_{20}^{\mathrm{pred}}-\rho_{20}|}{\rho_{20}}\approx 2.55\times 10^{-14}.
\]
However, it must be emphasized: as a power series at \(\theta=0\), its true convergence radius (measured in the unit-circle angle \(t\)) is pinned down by the nearest complex bifurcation point of Remark \ref{rem:sync-hatdelta-analytic-radius-complex} (\(R_t\approx 1.070953\)); therefore, the high precision at \(m=5\) (\(t=2\pi/5\)) should be understood as \emph{numerical extrapolation of a finite-order truncation}, rather than a convergent evaluation certifiable by geometric tail bounds; for auditable error budgets, one should follow the strategy fork of that remark.
\end{remark}

\begin{corollary}[\texorpdfstring{$m^{-2}$}{m^-2} gap law: closed-form constant for \texorpdfstring{$3-\rho_m$}{3-rho_m}]\label{cor:sync-rho-gap-m2}
Let \(\rho_m:=\max_{1\le r\le m-1}|\lambda(\omega_m^r)|\) as above, and set \(t=2\pi/m\).
Then as \(m\to\infty\) (equivalently \(t\to 0\)),
\[
\rho_m
=3\exp\!\Bigl(-\frac{11}{204}t^2+O(t^4)\Bigr)
=3-\frac{11}{68}t^2+O(t^4),
\]
yielding the explicit \(m^{-2}\) gap constant
\[
\boxed{\;
3-\rho_m\sim \frac{11\pi^2}{17}\cdot\frac{1}{m^2}\qquad(m\to\infty).
\;}
\]
If one further retains the fourth-order term, then
\[
\rho_m
=3\exp\!\Bigl(-\frac{11}{204}\Bigl(\frac{2\pi}{m}\Bigr)^2+\frac{1559}{1414944}\Bigl(\frac{2\pi}{m}\Bigr)^4+O(m^{-6})\Bigr),
\]
which already exhibits machine-precision agreement with the resultant true values of Table \ref{tab:sync_kernel_cyclotomic_elimination_summary} for moderate \(m\) (see the error comparison in Remark \ref{rem:sync-rho-pressure-asymp}).
\end{corollary}

\begin{corollary}[\texorpdfstring{$m^{-12}$}{m^-12} precision and \texorpdfstring{$m^2$}{m^2} diffusive scale]\label{cor:sync-rho-gap-m8-diffusive}
Continuing with the notation of Corollary \ref{cor:sync-rho-gap-m2}, and setting \(t=2\pi/m\).
From the \(\theta^{12}\) expansion of \(P(\theta)\) in Remark \ref{rem:sync-rho-pressure-asymp}, one obtains
\[
\begin{aligned}
\rho_m
&=3\exp\!\Bigl(-\frac{11}{204}t^2+\frac{1559}{1414944}t^4+\frac{17123893}{1177686190080}t^6\\
&-\frac{122803509253}{25412897733672960}t^8+\frac{518906628614669}{10575831520845339033600}t^{10}\\
&+\frac{24353138488976295223}{880247609142999258444595200}t^{12}+O(t^{14})\Bigr),
\end{aligned}
\]
whence
\[
\boxed{\;
\begin{aligned}
\rho_m
&=3-\frac{11\pi^2}{17m^2}
+\frac{1808\pi^4}{14739m^4}
-\frac{27872249\pi^6}{2044594080m^6}\\
&-\frac{77645433997\pi^8}{33089710590720m^8}
+\frac{2976802482983509\pi^{10}}{3442653489858508800m^{10}}\\
&+\frac{3832209193654464509\pi^{12}}{23878244605658617036800m^{12}}
+O(m^{-14}).\;
\end{aligned}\;}
\]
On the other hand, define the ``frequency cost/diffusion constant''
\[
\delta_m:=\log\frac{3}{\rho_m}=\log 3-\Re P\!\left(\frac{2\pi i}{m}\right),
\]
then
\[
\boxed{\;
\begin{aligned}
\delta_m
&=\frac{11\pi^2}{51m^2}
-\frac{1559\pi^4}{88434m^4}
-\frac{17123893\pi^6}{18401346720m^6}
+\frac{122803509253\pi^8}{99269131772160m^8}\\
&-\frac{518906628614669\pi^{10}}{10327960469575526400m^{10}}
-\frac{24353138488976295223\pi^{12}}{214904201450927553331200m^{12}}
+O(m^{-14}).\;
\end{aligned}\;}
\]
Therefore, for any residue-class error bound derived from the DFT/Witt interface
\(
\mathrm{err}_{n,m}\ll (\rho_m/3)^n
\)
(absorbing the constant into \(\ll\)), if one requires \(\mathrm{err}_{n,m}\le \varepsilon\), a sufficient condition is
\[
n\ \ge\ \frac{\log(1/\varepsilon)+O(1)}{\delta_m},
\]
yielding at leading order the diffusive time scale
\[
\boxed{\;
n_\varepsilon(m)\sim \frac{51}{11\pi^2}\,m^2\log\frac{1}{\varepsilon}\qquad(m\to\infty).\;}
\]
\end{corollary}

\paragraph{\texorpdfstring{$s$}{s}-plane zero-pole dictionary: spectral radius and ``no poles in the right half of the critical strip''}\label{par:s-plane-pole-dictionary}
The \(\zeta\)-function of this section, in the finite-matrix caliber, is a rational function, so all its singularities are completely determined by the finite spectrum.
To promote the ``spectral gap version of RH/GRH'' to the \((s)\)-plane language, introduce the Dirichlet coordinate
\[
z=\lambda^{-s},\qquad \lambda:=\rho(M)>0.
\]
\begin{remark}[Spectrum \(\Longleftrightarrow\) Dirichlet pole array]\label{prop:s-plane-spectrum-poles}
Let \(M\) be a finite-dimensional matrix, \(\lambda:=\rho(M)>0\). Define
\[
Z_M(s):=\frac{1}{\det\!\bigl(I-\lambda^{-s}M\bigr)}.
\]
Then:
\begin{enumerate}
  \item The poles of \(Z_M(s)\) are in one-to-one correspondence (counting algebraic multiplicity) with the nonzero eigenvalues \(\mu\) of \(M\). More specifically, \(\lambda^{-s}\mu=1\) if and only if
  \[
  s=\frac{\log\mu+2\pi i k}{\log\lambda}\qquad(k\in\ZZ),
  \]
  where \(\log\mu\) is any branch of the complex logarithm.
  \item The real part of the vertical line on which the poles lie is given by the modulus:
  \[
  \Re(s)=\frac{\log|\mu|}{\log\lambda}.
  \]
\end{enumerate}
Therefore, the kernel RH condition
\(
\max_{i\ge 2}|\mu_i|\le \sqrt{\lambda}
\)
is equivalent to: apart from the principal pole array (coming from \(\mu=\lambda\), corresponding to \(\Re(s)=1\)), all poles satisfy \(\Re(s)\le \tfrac12\), i.e., the region \(\Re(s)>\tfrac12\) contains no nontrivial pole contributions.
\end{remark}

\paragraph{Finite-dimensional explicit formula: the error in prime shadows is the ``critical-circle spectral contribution''}\label{par:finite-explicit-formula}
For finite-dimensional kernels, the ``leading term + error'' of primitive orbit counting can be written entirely as a finite sum over eigenvalues, thereby directly identifying the square-root oscillation as the spectral contribution from the critical circle \(|\mu|=\sqrt{\lambda}\).
\begin{remark}[Finite-dimensional explicit formula template]\label{prop:finite-dimensional-explicit-formula}
Let the eigenvalues of \(M\) (counted with algebraic multiplicity) be \(\mu_1,\dots,\mu_d\), and write
\(
a_n:=\Tr(M^n)=\sum_{j=1}^d \mu_j^n
\).
Let \(p_n\) denote the number of primitive orbits of length \(n\) (counted by orbits); then M\"obius inversion gives
\[
p_n=\frac{1}{n}\sum_{m\mid n}\mu(m)\,a_{n/m}
=\frac{1}{n}\sum_{j=1}^d\ \sum_{m\mid n}\mu(m)\,\mu_j^{\,n/m}.
\]
In particular, if \(\lambda:=\rho(M)\) is a simple Perron root, then
\[
p_n=\frac{\lambda^n}{n}+O\!\left(\frac{\Lambda^n}{n}\right),
\qquad
\Lambda:=\max_{j\ge 2}|\mu_j|.
\]
If there exists a nontrivial eigenvalue satisfying \(|\mu|=\sqrt{\lambda}\), then in general an non-improvable square-root-level oscillation term (\(\Omega(\lambda^{n/2}/n)\)) appears.
\end{remark}

\begin{proposition}[RH stratification for three classes of kernels: synchronization kernel non-RH, real-input kernel RH-sharp, low-order collision kernels strict RH]\label{prop:rh-stratification-three-kernels}
In the unweighted case ($u=1$), the three typical classes of finite kernels appearing in this paper satisfy the following auditable stratification:
\begin{enumerate}
  \item \textbf{Synchronization kernel: non-RH.} The 10-state synchronization kernel $B$ does not satisfy the kernel RH: there exists a nontrivial eigenvalue with \(|\mu|>\sqrt{\lambda}\) (Remark \ref{prop:sync-kernel-10-nonrh}).
  \item \textbf{Real-input kernel: RH-sharp.} The kernel RH for the 40-state real-input kernel $M$ is achieved with equality: there exists a boundary-attaining eigenvalue \(\mu=-\sqrt{\lambda}\), so that the square-root-level oscillation cannot be eliminated (Proposition \ref{prop:finite-rh-40} and Remark \ref{prop:finite-dimensional-explicit-formula}).
  \item \textbf{Folding collision kernels: strict RH.} The non-principal spectra of the second-order and third-order collision kernels \(A_2,A_3\) lie strictly inside the square-root disk, i.e.,
  \[
  \max_{i\ge 2}|\lambda_i(A_q)|<\sqrt{\rho(A_q)}\qquad(q=2,3),
  \]
  so their prime-orbit error scale is strictly better than the square-root upper bound (Remarks \ref{rem:pom-collision-rh-margin} and \ref{rem:pom-collision-rh-margin-a3}).
\end{enumerate}
\end{proposition}

\begin{remark}[Near-RH violation by higher-order folding collision kernels: starting from \texorpdfstring{$q=4$}{q=4}]\label{rem:collision-kernel-rh-break-q4}
The ``strict RH for folding collision kernels'' stated above applies only to the low-order cases \(A_2,A_3\).
The fourth-order five-state kernel \(A_4\) already exhibits a boundary-level violation, and within the verified range \(q\le 8\), the violation for \(q\ge 5\) increases rapidly (Remark \ref{rem:pom-collision-rh-break-a4}, Table \ref{tab:fold_collision_kernel_rh_scan_q2_8}). The boundary-level violation of \(A_4\) can be further attributed to a single-parameter edge weight critical value \(t_\star\) (Remark \ref{rem:pom-collision-rh-break-a4-threshold}).
\end{remark}

\begin{corollary}[\texorpdfstring{$s$}{s}-plane geometric discrimination: pole hierarchy in the right half of the critical strip]\label{cor:rh-stratification-splane}
Let \(Z_M(s)=\det(I-\lambda^{-s}M)^{-1}\) be the Dirichlet-transformed \(\zeta\)-function (Remark \ref{prop:s-plane-spectrum-poles}). Then the RH stratification of the above three classes of kernels can be equivalently translated into the pole geometric hierarchy in the \(s\)-plane:
\begin{enumerate}
  \item \textbf{Synchronization kernel: non-RH.} Due to the existence of a nontrivial eigenvalue with \(|\mu|>\sqrt{\lambda}\), there exists a sequence of poles lying on the vertical line
  \(
  \Re(s)=\log|\mu|/\log\lambda
  \),
  and this real part is strictly greater than \(1/2\). Therefore, nontrivial poles appear in the right half-open strip \(\tfrac12<\Re(s)<1\).
  \item \textbf{Real-input kernel: RH-sharp.} All nontrivial pole arrays in the critical strip lie exactly on \(\Re(s)=1/2\) (Proposition \ref{prop:finite-rh-40} of Appendix \ref{app:real-input-40-finite-rh}).
  \item \textbf{Folding collision kernels: strict RH.} For \(A_2,A_3\), one has \(\max_{i\ge 2}|\mu_i|<\sqrt{\lambda}\), so the region \(\tfrac12\le \Re(s)<1\) contains no nontrivial poles; in other words, there are no boundary-attaining pole array contributions on the critical line \(\Re(s)=1/2\) either (apart from the principal pole array at \(\Re(s)=1\)).
\end{enumerate}
\end{corollary}
\begin{proof}
By Remark \ref{prop:s-plane-spectrum-poles}, each nonzero eigenvalue \(\mu\) generates a pole array whose vertical line has real part \(\Re(s)=\log|\mu|/\log\lambda\). The three cases correspond to:
non-RH when \(|\mu|>\sqrt{\lambda}\Rightarrow \Re(s)>1/2\);
RH-sharp when \(|\mu|=\sqrt{\lambda}\Rightarrow \Re(s)=1/2\);
strict RH when for all nontrivial \(|\mu|<\sqrt{\lambda}\Rightarrow \Re(s)<1/2\).
\qed
\end{proof}

\paragraph{Spectral reciprocity criterion for the \texorpdfstring{$s$}{s}-plane completed functional equation}\label{par:s-plane-functional-eq-criterion}
The ``functional equation'' in the preceding paragraph takes place on the parameter side ($u\leftrightarrow 1/u$ or $\theta\leftrightarrow-\theta$).
If one wishes to further promote the symmetry to an
\(
s\leftrightarrow 1-s
\)
type symmetry in the Dirichlet coordinate $z=\lambda^{-s}$, one requires a spectral reciprocity structure independent of the parameter self-duality.

\begin{proposition}[\texorpdfstring{$\lambda$}{lambda}-reciprocity (spectral reciprocal closure) \texorpdfstring{$\Longleftrightarrow$}{<->} completed functional equation]\label{prop:reciprocal-spectrum-functional-equation}
Let $M$ be a finite-dimensional matrix, $\lambda:=\rho(M)>0$ its Perron root. Set
\[
\Delta_M(z):=\det(I-zM),\qquad
Z_M(s):=\Delta_M(\lambda^{-s})^{-1},
\qquad
r:=\deg_z\Delta_M(z).
\]
and define the completed function
\[
\Xi_M(s):=\lambda^{-rs/2}\,Z_M(s).
\]
Then the following three statements are equivalent:
\begin{enumerate}
  \item \textbf{Spectral reciprocal closure:} For every \(\mu\in\mathrm{spec}(M)\setminus\{0\}\),
  \[
  \frac{\lambda}{\mu}\in\mathrm{spec}(M)\setminus\{0\}
  \qquad(\text{counted with algebraic multiplicity}).
  \]
  \item \textbf{\(\lambda\)-reciprocity of \(\Delta_M\):} There exists a constant \(C'\neq 0\) such that for all complex numbers \(z\),
  \[
  \lambda^{r}z^{r}\Delta_M\!\left(\frac{1}{\lambda z}\right)=C'\Delta_M(z).
  \]
  \item \textbf{Completed functional equation:} There exists a constant \(C\neq 0\) such that
  \[
  \Xi_M(s)=C\,\Xi_M(1-s),
  \qquad C=C'\lambda^{-r/2}.
  \]
\end{enumerate}
\end{proposition}
\begin{proof}
First prove (1)\(\Rightarrow\)(2). Denote all nonzero eigenvalues of \(M\) (counted with algebraic multiplicity) by \(\mu_1,\dots,\mu_r\); then
\[
\Delta_M(z)=\prod_{j=1}^{r}(1-\mu_j z).
\]
Therefore
\[
\begin{aligned}
\lambda^r z^r\Delta_M\!\left(\frac{1}{\lambda z}\right)
&=\lambda^r z^r\prod_{j=1}^r\left(1-\frac{\mu_j}{\lambda z}\right)
\\
&=\lambda^r z^r\prod_{j=1}^r\left(-\frac{\mu_j}{\lambda z}\right)\prod_{j=1}^r\left(1-\frac{\lambda}{\mu_j}z\right)
\\
&=(-1)^r\Bigl(\prod_{j=1}^r\mu_j\Bigr)\prod_{j=1}^r\left(1-\frac{\lambda}{\mu_j}z\right).
\end{aligned}
\]
By the reciprocal closure of (1), \(\{\lambda/\mu_j\}_{j=1}^r\) and \(\{\mu_j\}_{j=1}^r\) are the same as multisets, so
\[
\prod_{j=1}^r\left(1-\frac{\lambda}{\mu_j}z\right)=\prod_{j=1}^r(1-\mu_j z)=\Delta_M(z).
\]
Therefore (2) holds, with \(C'=(-1)^r\prod_{j=1}^r\mu_j\neq 0\).
\par
Next prove (2)\(\Rightarrow\)(3). Substitute \(z=\lambda^{-s}\) in (2), noting that
\(
(\lambda z)^{-1}=\lambda^{-(1-s)}
\), to obtain
\[
\lambda^{r(1-s)}\,\Delta_M\!\bigl(\lambda^{-(1-s)}\bigr)=C'\,\Delta_M(\lambda^{-s}).
\]
Taking reciprocals on both sides and rearranging gives
\[
Z_M(1-s)=\Delta_M\!\bigl(\lambda^{-(1-s)}\bigr)^{-1}=\frac{\lambda^{r(1-s)}}{C'}\,\Delta_M(\lambda^{-s})^{-1}=\frac{\lambda^{r(1-s)}}{C'}\,Z_M(s).
\]
Then
\[
\Xi_M(1-s)
=\lambda^{-r(1-s)/2}Z_M(1-s)
=\frac{\lambda^{r(1-s)/2}}{C'}\,Z_M(s),
\]
while
\(
\Xi_M(s)=\lambda^{-rs/2}Z_M(s)
\).
Hence
\(
\Xi_M(s)=C'\lambda^{-r/2}\,\Xi_M(1-s)
\), i.e., (3) holds with \(C=C'\lambda^{-r/2}\).
\par
Finally prove (3)\(\Rightarrow\)(1). Since \(\Xi_M(s)=\lambda^{-rs/2}Z_M(s)\), both have the same pole set and pole orders.
If \(s_0\) is a pole of \(Z_M\), then by (3), \(1-s_0\) is also a pole with the same order.
By the ``spectrum \(\leftrightarrow\) pole'' dictionary of Remark \ref{prop:s-plane-spectrum-poles},
the pole \(s_0\) corresponds to some \(\mu\in\mathrm{spec}(M)\setminus\{0\}\) with \(\lambda^{-s_0}\mu=1\), so \(\lambda^{-(1-s_0)}=\mu/\lambda\).
That \(1-s_0\) is a pole means there exists \(\mu'\in\mathrm{spec}(M)\setminus\{0\}\) with \(\lambda^{-(1-s_0)}\mu'=1\), i.e., \((\mu/\lambda)\mu'=1\), so \(\mu'=\lambda/\mu\).
The equality of pole orders gives the algebraic multiplicity consistency, so (1) holds.
\end{proof}

\begin{remark}[Equivalent \texorpdfstring{$\Delta_M$}{Delta}-reciprocal polynomial criterion]\label{rem:delta-reciprocal-polynomial}
The spectral condition of Proposition \ref{prop:reciprocal-spectrum-functional-equation} is equivalent to the $z$-root set being invariant (with multiplicity) under
\(
z\mapsto (\lambda z)^{-1}
\),
which is also equivalent to the existence of $C'\neq 0$ such that
\[
\lambda^{r}z^{r}\Delta_M\!\left(\frac{1}{\lambda z}\right)=C'\Delta_M(z).
\]
Therefore, when the coefficients of $\Delta_M$ are exactly computable (e.g., for integer-coefficient matrices), the above formula provides a term-by-term auditable criterion.
\end{remark}

\begin{proposition}[Reciprocal closure \texorpdfstring{$+$}{+} kernel RH \texorpdfstring{$\Rightarrow$}{=>} toy-RH]\label{prop:reciprocal-plus-kernel-rh-implies-toy-rh}
Let \(M\) be a finite-dimensional matrix, \(\lambda=\rho(M)>1\), and set
\(
Z_M(s)=\det(I-\lambda^{-s}M)^{-1}
\).
If \(M\) satisfies the reciprocal closure of Proposition \ref{prop:reciprocal-spectrum-functional-equation} (equivalently: \(\Delta_M\) is a \(\lambda\)-reciprocal polynomial), and also satisfies the kernel RH (Ramanujan radius bound)
\[
\max_{\mu\in\mathrm{spec}(M),\,\mu\neq \lambda}|\mu|\le \sqrt{\lambda},
\]
then \(Z_M\) satisfies toy-RH: all nontrivial poles in the open strip \(0<\Re(s)<1\) lie on \(\Re(s)=\tfrac12\).
Equivalently,
\[
(\forall\,\mu\ \text{satisfying}\ 1<|\mu|<\lambda)\qquad |\mu|=\sqrt{\lambda}.
\]
Thus, when the reciprocal closure holds, the kernel RH and toy-RH are equivalent.
\end{proposition}
\begin{proof}
Let \(\mu\in\mathrm{spec}(M)\) satisfy \(1<|\mu|<\lambda\). By the kernel RH, \(|\mu|\le \sqrt{\lambda}\).
If \(|\mu|<\sqrt{\lambda}\), then by the reciprocal closure \(\lambda/\mu\in\mathrm{spec}(M)\), and
\(
|\lambda/\mu|>\sqrt{\lambda}
\), contradicting the kernel RH.
Hence for all eigenvalues satisfying \(1<|\mu|<\lambda\), one must have \(|\mu|=\sqrt{\lambda}\), which is equivalent to toy-RH (by the dictionary of Remark \ref{prop:s-plane-spectrum-poles}, \(\Re(s)=\log|\mu|/\log\lambda\)).
\par
Conversely, toy-RH implies that no eigenvalue satisfies \(\sqrt{\lambda}<|\mu|<\lambda\), since such an eigenvalue would produce an off-critical-line pole with \(\tfrac12<\Re(s)<1\); and \(|\mu|\le \lambda\) always holds (\(\lambda=\rho(M)\)), so toy-RH \(\Rightarrow\) kernel RH.
Therefore, under the reciprocal closure condition, the two are equivalent.
\end{proof}

\begin{proposition}[Vertical lattice obstruction for finite-dimensional determinants]\label{prop:finite-kernel-lattice-obstruction}
Let \(M\) be a finite-dimensional matrix, \(\lambda:=\rho(M)>1\), and define
\[
Z_M(s):=\det(I-\lambda^{-s}M)^{-1}.
\]
Then the set of all poles of \(Z_M(s)\) (counted with multiplicity) forms the finite union
\[
\mathcal{S}(M)=\bigcup_{\mu\in\mathrm{spec}(M)\setminus\{0\}}\ \Bigl\{\frac{\log\mu+2\pi i k}{\log\lambda}\;:\;k\in\ZZ\Bigr\},
\]
where \(\log\mu\) is taken over any complex logarithm branch. In particular, each nonzero eigenvalue \(\mu\) contributes a \textbf{strictly arithmetic} vertical pole array with imaginary part spacing
\[
\Delta=\frac{2\pi}{\log\lambda}.
\]
More generally, take any finite collection of triples \((M_j,\lambda_j,\varepsilon_j)\) (where \(M_j\) is finite-dimensional, \(\lambda_j>1\), \(\varepsilon_j\in\{\pm 1\}\)), and consider the finite product/quotient
\[
F(s):=\prod_{j=1}^{J}\det(I-\lambda_j^{-s}M_j)^{-\varepsilon_j},
\]
then the zero and pole set of \(F\) is a \textbf{finite union of finitely many vertical arithmetic progressions}, so its counting function satisfies
\[
N_F(T):=\#\{\,s\in\mathrm{Zeros}(F)\cup\mathrm{Poles}(F):\ 0<\Im(s)\le T\,\}=O(T).
\]
Therefore, no fixed finite product \(F\) can have its ``zero/pole set in the critical strip'' coincide with the nontrivial zero set of the Riemann \(\zeta(s)\) (or \(\xi(s)\)): the latter satisfies \(N_\zeta(T)\asymp T\log T\) by the Riemann--von Mangoldt formula.
\end{proposition}
\begin{proof}
The first paragraph is a direct expansion of Remark \ref{prop:s-plane-spectrum-poles}: \(\det(I-\lambda^{-s}M)=0\) if and only if \(\lambda^{-s}=1/\mu\) (for some \(\mu\neq 0\)); taking logarithms yields the pole array formula; the imaginary part spacing comes from the \(2\pi i\)-multivaluedness of \(\log(\cdot)\).
\par
For the finite product \(F\), its zeros/poles can only come from the union of the zero sets of the individual factors \(\det(I-\lambda_j^{-s}M_j)\); each factor contributes \(O(T)\) arithmetic points in \((0,T]\) (with constants depending only on \(\log\lambda_j\) and \(\#(\mathrm{spec}(M_j)\setminus\{0\})\)), so the union is still \(O(T)\).
\par
The final paragraph: the nontrivial zero counting of \(\zeta(s)\) satisfies \(N_\zeta(T)=\frac{T}{2\pi}\log\frac{T}{2\pi}-\frac{T}{2\pi}+O(\log T)\), so \(N_\zeta(T)\) is incompatible in growth order with \(O(T)\), and thus cannot coincide with the zero/pole set of any fixed finite product \(F\). \qed
\end{proof}

\begin{corollary}[Arithmetic lattice closure: the imaginary part difference set is constrained by a finitely generated additive subgroup]\label{cor:finite-determinant-arithmetic-lattice}
Continuing with the notation of Proposition \ref{prop:finite-kernel-lattice-obstruction}, and let \(F\) be the finite product/quotient therein.
Then there exist a finite set \(\Sigma\subset\CC\) and a finite set \(\{\Delta_1,\dots,\Delta_J\}\subset(0,\infty)\) (where each \(\Delta_j=2\pi/\log\lambda_j\)) such that
\[
\mathrm{Zeros}(F)\cup\mathrm{Poles}(F)
=\bigcup_{\ell}\Bigl\{\ \sigma_\ell+i(t_\ell+k\Delta_{j(\ell)})\ :\ k\in\ZZ\ \Bigr\}
\qquad(\sigma_\ell\in\RR,\ t_\ell\in\RR).
\]
Therefore, for any real-part strip \(a<\Re(s)<b\), let
\[
\mathcal{T}_{F}(a,b)
:=\bigl\{\Im(s):\ s\in\mathrm{Zeros}(F)\cup\mathrm{Poles}(F),\ a<\Re(s)<b\bigr\}\subset\RR,
\]
then its difference set satisfies
\[
\mathcal{T}_{F}(a,b)-\mathcal{T}_{F}(a,b)
\subseteq
\mathcal{E}\ +\ \sum_{j=1}^{J}\Delta_j\ZZ,
\]
where \(\mathcal{E}\subset\RR\) is some finite set. In particular, for each spacing \(\Delta_j\) that actually appears in the strip, the difference set contains a strict arithmetic sublattice \(\Delta_j\ZZ\).
\end{corollary}
\begin{proof}
Proposition \ref{prop:finite-kernel-lattice-obstruction} has already established that the zero/pole set is a finite union of finitely many vertical arithmetic progressions; writing each point on each vertical line in the form \(\sigma_\ell+i(t_\ell+k\Delta_{j(\ell)})\) yields the first formula.
\par
For the difference set, take any two points \(t=t_\ell+k\Delta_{j(\ell)}\), \(t'=t_{\ell'}+k'\Delta_{j(\ell')}\); then
\[
t-t'=(t_\ell-t_{\ell'})+\bigl(k\Delta_{j(\ell)}-k'\Delta_{j(\ell')}\bigr)\in (t_\ell-t_{\ell'})+\sum_{j=1}^J\Delta_j\ZZ,
\]
and the finite union introduces only finitely many offset sets
\(
\mathcal{E}:=\{t_\ell-t_{\ell'}\}_{\ell,\ell'}.
\)
The final statement: if some arithmetic progression with spacing \(\Delta_j\) appears in the strip, then that progression contains infinitely many points in the imaginary direction, so its internal difference set must contain \(\Delta_j\ZZ\). \qed
\end{proof}

\begin{corollary}[Statistical lattice obstruction: the two-point difference must contain Dirac combs]\label{cor:finite-determinant-statistical-lattice}
Under the conditions of Corollary \ref{cor:finite-determinant-arithmetic-lattice}, if \(\mathcal{T}_F(a,b)\) contains a complete arithmetic progression with spacing \(\Delta>0\), namely \(t_0+\Delta\ZZ\) (e.g., if an eigenvalue array of some factor \(\det(I-\lambda^{-s}M)\) has not been completely canceled), then define the two-point difference measure
\[
\nu_T:=\frac{1}{T}\sum_{t,t'\in\mathcal{T}_F(a,b)\cap(0,T]}\delta_{t-t'}\qquad(T>0),
\]
then for every \(n\in\ZZ\),
\[
\liminf_{T\to\infty}\ \nu_T\bigl(\{n\Delta\}\bigr)\ \ge\ \frac{1}{\Delta}\ >0.
\]
Therefore, any limit two-point difference statistic as \(T\to\infty\) must contain an atomic part supported on \(\Delta\ZZ\); this rules out all target statistical forms in which ``the two-point difference after unfolding has a continuous density.''
\end{corollary}
\begin{proof}
Let \(\mathcal{A}_T:=(t_0+\Delta\ZZ)\cap(0,T]\), so \(|\mathcal{A}_T|=T/\Delta+O(1)\).
For any integer \(n\), the number of ordered pairs with difference \(t-t'=n\Delta\) includes at least \(\{(t_0+(k+n)\Delta,\ t_0+k\Delta)\}\) (when both points fall in \((0,T]\)), so
\[
\nu_T(\{n\Delta\})
\ge
\frac{1}{T}\#\bigl\{(t,t')\in\mathcal{A}_T^2:\ t-t'=n\Delta\bigr\}
=\frac{|\mathcal{A}_T|-O(1)}{T}
=\frac{1}{\Delta}+o(1).
\]
Taking \(\liminf\) yields the conclusion. \qed
\end{proof}

\begin{proposition}[Critical-circle mode count \texorpdfstring{$B$}{B} and the linear calibration of critical-line zero-pole density]\label{prop:critical-circle-density-calibration}
Let \(M\) be a finite-dimensional matrix, \(\lambda=\rho(M)>1\), and let
\(
Z_M(s)=\det(I-\lambda^{-s}M)^{-1}
\).
Denote all boundary-attaining eigenvalues on the critical circle (counted with algebraic multiplicity) by
\[
\mu_1,\dots,\mu_B,\qquad |\mu_b|=\sqrt{\lambda}.
\]
Then all poles of \(Z_M\) on the critical line \(\Re(s)=\tfrac12\) are given by \(B\) vertical arithmetic progressions:
\[
s=\frac12+\frac{i(\arg\mu_b+2\pi k)}{\log\lambda},\qquad b=1,\dots,B,\ k\in\ZZ,
\]
with imaginary part spacing \(\Delta=2\pi/\log\lambda\). Moreover, the critical-line pole count up to height \(T\) satisfies
\[
N_{\mathrm{crit}}(T)
:=\sum_{\substack{s_0\in\mathrm{Poles}(Z_M)\\ \Re(s_0)=\tfrac12,\ 0<\Im(s_0)\le T}}\mathrm{ord}_{s_0}Z_M
=\frac{B\log\lambda}{2\pi}\,T+O(1)\qquad(T\to\infty).
\]
\end{proposition}
\begin{proof}
By the pole dictionary of Proposition \ref{prop:finite-kernel-lattice-obstruction}, \(\mu_b\) produces the pole array
\(
s=(\log\mu_b+2\pi ik)/\log\lambda
\).
When \(|\mu_b|=\sqrt{\lambda}\), one has \(\Re(s)=\log|\mu_b|/\log\lambda=\tfrac12\), and writing \(\log\mu_b\) as \(\tfrac12\log\lambda+i\arg\mu_b\) gives the critical-line parametrization; the spacing is given by \(2\pi i/\log\lambda\).
\par
For the counting, each arithmetic progression contributes \(T/\Delta+O(1)=(\log\lambda/(2\pi))T+O(1)\) points in \((0,T]\); summing over \(b=1,\dots,B\) gives the stated leading term and \(O(1)\) remainder. \qed
\end{proof}

\begin{corollary}[Necessary condition for \texorpdfstring{$T\log T$}{T log T} density matching: \texorpdfstring{$B(T)\asymp (\log T)/(\log\lambda)$}{B(T) ~ (log T)/(log lambda)}]\label{cor:critical-circle-mode-growth}
If one wishes, within the ``kernel RH/RH-sharp critical-line mode'' framework, to promote the critical-line zero-pole density up to height \(T\) to the \(N(T)\asymp T\log T\) level (as in the Riemann--von Mangoldt leading term), then for any fixed \(\lambda>1\) finite-dimensional kernel, it is necessary to introduce a family of kernels \(M_T\) growing with scale such that their critical-circle mode count \(B=B(T)\) satisfies
\[
B(T)\ \asymp\ \frac{\log T}{\log\lambda}.
\]
For the 40-state real-input kernel with \(\lambda=\varphi^2\), one has \(B(T)\asymp (\log T)/(2\log\varphi)\).
\end{corollary}
\begin{proof}
By Proposition \ref{prop:critical-circle-density-calibration}, for fixed \(\lambda\), the critical-line point count satisfies \(N_{\mathrm{crit}}(T)=\frac{B\log\lambda}{2\pi}T+O(1)\).
To achieve \(N_{\mathrm{crit}}(T)\asymp T\log T\), one must have \(B=B(T)\asymp (\log T)/(\log\lambda)\).
Substituting \(\lambda=\varphi^2\) and \(\log\lambda=2\log\varphi\) yields the result. \qed
\end{proof}

\begin{remark}[Explicit lift generator for \texorpdfstring{$\lambda=\varphi^2$}{lambda=phi^2}: using cyclic lifts to make \texorpdfstring{$B$}{B} grow linearly]\label{rem:critical-circle-mode-growth-cyclic-lift}
Appendix \ref{app:real-input-40-finite-rh} has already provided an auditable lift generator: take \(q\ge 2\), let \(P_q\) be the \(q\)-cyclic permutation matrix and define the cyclic lift
\(
M^{\langle q\rangle}:=M\otimes P_q
\) (Proposition \ref{prop:finite-rh-phase-lift}).
For the 40-state real-input kernel, the unique boundary-attaining eigenvalue on the critical circle is \(-\sqrt{\lambda}=-\varphi\) (Remark \ref{rem:finite-rh-40-ramanujan}), so the cyclic lift amplifies the critical-circle mode count to \(B=q\) and refines the critical-line comb spacing to
\(
2\pi/(q\log\lambda)
\) (Remark \ref{rem:finite-rh-height-scale-calibration}).
Therefore, if \(q\) serves as the lift fiber size, the necessary scale law for density matching can be written as
\[
q(T)\ \asymp\ \frac{\log T}{\log\lambda}
=\frac{\log T}{2\log\varphi},
\]
which constitutes a fully explicit instance of the ``lift--scale calibration law.''
\end{remark}

\begin{remark}[Verification outcome for the two classes of kernels in this section: reproducible]\label{rem:reciprocal-check-outcome}
The $s$-plane symmetry described above does not hold automatically: it requires the nonzero spectrum to be closed under $\mu\mapsto\lambda/\mu$.
For the 10-state synchronization kernel $B$, from the factor \((z+1)\) of \(\det(I-zB)\) one knows \(\mu=-1\in\mathrm{spec}(B)\), but \(\lambda/\mu=-3\notin\mathrm{spec}(B)\), so the reciprocal closure fails.
For the 40-state real-input kernel $M$, Theorem \ref{thm:pom-rh-sharp} gives \(\lambda^{-1}\in\mathrm{spec}(M)\setminus\{0\}\), but \(\lambda/(\lambda^{-1})=\lambda^2\notin\mathrm{spec}(M)\setminus\{0\}\), so it also fails.
The corresponding spectral reciprocal closure/coefficient reciprocity tests are automatically output by the script \texttt{scripts/exp\_reciprocal\_spectrum\_functional\_equation\_check.py} to
\texttt{artifacts/export/reciprocal\_spectrum\_functional\_equation\_check.json}.
\end{remark}
